\section{Extra decay from the degree of a cluster}
\label{sec:high-degree-cluster-prices}
After extracting the singleton reference, many terms can be paid
without a further subtraction. The relevant degree counts elementary
occurrences inside every connected component, together with an extra
grade for certain species. Counting only the number of components
in a cluster would miss this gain. We follow the two successive
connected expansions and retain their separate convergence conditions.

\subsection{Grading the elementary activities}
Write $\beta_+(x)=\beta_F(G,x,L)$ for the shifted fine coupling and set
\[
                         D(x)=\beta_+(x)^{3/8}.
\]
Eventually $D\ge1$. Give the six inherited logarithmic species the
following additional grades:
\[
\begin{array}{c|rrrrrr}
\text{species}&\text{regular}&\text{cubic}&\text{quartic}
 &\text{quadratic loop}&\text{fine tail}&\text{coarse tail}\\\hline
 h&1&0&1&1&2&2.
\end{array}
\]
These are extra grades; a further factor $D$ per occurrence will
be inserted at the connected-component stage.

Section~\ref{sec:inherited-activity-prices} supplies the required
full-support site bounds. Fix $-1/2<t<0$, and choose $s$ strictly
between $-1/2$ and $t$. The ungraded regular, quartic and loop prices
can each be bounded by $\beta_+^{-15/16}$ eventually. Multiplication
by $D$ changes this exponent to $-9/16<s$. The cubic price has upper
exponent $-1/2$ and is bounded by $\beta_+^s$. Both true path tails
can be bounded by $\beta_+^{-2}$; multiplication by $D^2$ changes this
to $-5/4<s$. The six prices therefore sum to at most
$6\beta_+^s\le\beta_+^t$ eventually. These inequalities preserve the
original supports and all length charges. The forced inherited family
retains every negative upper exponent and receives no extra grade.

The same argument applies to the local derivatives. Hold the original
insertion fixed in the compensated coordinates as in
Section~\ref{sec:inherited-activity-prices}. Its transported cutoff has
jet scale bounded by an upper-exponent-zero function times the full
cover plus length charge. The coarse tail has an additional varying
coarse-field term; its own jet price has upper exponent $-2$, and so
the factor $D^2$ again leaves exponent $-5/4$. Adding this term to its
old price preserves the common bound $\beta_+^t$ by first taking the
intermediate exponent $s$. The cutoff product rule only changes the
upper-exponent-zero jet scale.

For the other optional letters, assign grade zero to the quadratic
local term and grade two to the remaining small correction family.
Their ungraded prices are, respectively, of upper exponent $-1/2$
and at most $-2$. The density and exponential jet estimates of
Section~\ref{sec:activity-jet-assembly} therefore preserve a common
upper exponent $t$ after the grades are inserted. More explicitly,
if $a_i$ is a nonnegative logarithmic amplitude and its jet scale is
$b_i$, the full-support site bound implies $a_i\le1$ eventually:
choose a root in its nonempty support, bound this single nonnegative
summand by the site sum, and drop its weights, all at least one.
The zero-order exponential difference is bounded by $e a_i$.
For a positive-order derivative, Fa\`a di Bruno's formula is a finite
sum indexed by partitions; a partition with $q$ blocks contributes
at most $e a_i^q b_i^r$ at order $r$. Since $a_i^q\le a_i$, this
bounds every jet through order $n$ by
$C_n a_i\max(1,b_i)^n$. The scale $b_i$ retains its full cover and
length charges; only its scalar prefactor has upper exponent zero.
The polynomial charge factor is paid by the reserve in the next
paragraph.
The previously proved warning
estimates still give the same joint warning power after multiplying
by these deterministic grades. Thus the grading is compatible with
the actual Gaussian derivative majorants, not only their zeroth-order
values.

The passage from graded letter prices to differentiated component
integrals uses the common random field. For a component $P$, take
$I$ to be the full set of fine bonds read by its enlarged union
support, and let $J$ be its joint warning event. For the Gaussian
coordinates $z_{b,c}$ indexed by bond and colour, define
\[
 \mathcal F_I(z)=1+|I|+|I\times\mathcal C|
                   +\sum_{(b,c)\in I\times\mathcal C}|z_{b,c}|,
\]
where $\mathcal C$ is the finite colour index set.
The product-rule
bound has the form
\[
 A_P\mathbf1_J\mathcal F_I^{kj},
 \qquad
 A_P=\prod_{i\in P}\alpha_i\,V(x)(1+S(P))^j,
\]
where $\alpha_i$ are the nonnegative elementary packet and letter prices,
$V\ge0$ has upper exponent zero, and $S(P)$ is the full footprint
charge. The read-count bound is
$|I|\le(1+S(P))b_B d_L$, with the fixed block-density factor $d_L$.
The Gaussian moment estimate of
Section~\ref{sec:gaussian-moment-prices}, at the doubled order $2kj$,
gives
$\int\mathcal F_I^{2kj}\,d\mu\le
[E(x)(1+S(P))^K]^2$
for some fixed $K$ and upper-exponent-zero $E\ge0$.
The joint warning bound is
$\mu(J)\le\rho^{2b(P)}$, where $b(P)$ is the number of warning
labels and $\rho\ge0$ is the retained warning price. Cauchy--Schwarz
therefore yields
\[
 \int A_P\mathbf1_J\mathcal F_I^{kj}\,d\mu
 \le A_PE(x)(1+S(P))^K\rho^{b(P)}.
\]
The remaining Haar integral has mass one. This is precisely the
integrated component majorant, with its same warning power and a
larger fixed footprint exponent. It uses no independence between
overlapping reads or warning events. The local jet construction and
its common integrable majorant justify differentiation as proved in
Section~\ref{sec:activity-jet-assembly}.

\subsection{Recovering the footprint weight in the first expansion}
We state the needed positive-weight calculation explicitly. For an
elementary label $i$, let $X_i$ be its support, $u_i\ge0$ its price,
and $e_i\ge0$ its additional length charge. For a finite connected
set $P$ put
\[
 C(P)=c\left(\bigcup_{i\in P}X_i\right),\qquad
 S(P)=\sum_{i\in P}(c(X_i)+e_i).
\]
Here $c$ is the connected cover number at the current scale. The
cover subadditivity along a spanning tree gives $C(P)\le\sum_i c(X_i)$.
Fix $a\ge0$, $\epsilon>0$ and integer $k\ge0$. If the elementary
site estimate is available with exponent $a+2\epsilon$ on cover and
$2\epsilon$ on the length charge, the first tree estimate yields
its rooted connected sum. It also pays the stronger summand
\[
 \mathbf1_{\{z\in\bigcup X_i\}}
       \prod_{i\in P}u_i\,e^{aC(P)}
                   (1+S(P))^k e^{\epsilon S(P)}.
\]
Indeed
\[
 (1+S)^k\le e^\epsilon k!\epsilon^{-k}e^{\epsilon S},
 \qquad S\ge0,
\]
by the exponential series at $\epsilon(1+S)$. Together with cover
subadditivity this bounds the displayed summand by
$e^\epsilon k!\epsilon^{-k}$ times the summand already paid by the
elementary site estimate. The same $\epsilon$ therefore pays both
the retained full-footprint exponential and the polynomial moment;
we have used two copies of its reserve, not discarded either cost.

Apply this calculation to the mixed elementary alphabet of
Section~\ref{sec:mixed-activity-prices}. Its mandatory packet weights
already include their seed payment, while its optional letters have
the grades above. For a component $P$, let $h(P)$ be the sum of its
extra grades, $|P|$ its number of elementary occurrences, $m(P)$ its
number of mandatory seeds, and $S(P)$ its full footprint charge.
The strengthened component multiplier is
\begin{equation}\label{eq:degree-component-multiplier}
 D^{h(P)+|P|}\exp\{M m(P)+\epsilon S(P)\},
\end{equation}
where $M\ge0$ is the seed mark. This factor is the product of the
elementary factors $D^{1+h_i}e^{M m_i}$ and the retained footprint
factor. Multiplying every elementary price by $D$ multiplies its
site budget by $D$. Consequently the first expansion is valid if
\begin{equation}\label{eq:degree-first-smallness}
 \frac{4b_B D r_0(x)}{\delta_0}\le\frac12,
\end{equation}
where $r_0$ and $\delta_0>0$ are the original elementary root price
and its reserve, and $b_B=(2B+1)^4$. Its component root price is
$D\omega(x)$, where $\omega$ is the original component price including
the footprint polynomial just proved. All these statements require
the graded elementary site bound on the same incoming representation.

\subsection{The second expansion and its degree gain}
Let $j$ be the field derivative order, and let $B_1(x)\ge0$ be the
upper-exponent-zero constant from the component derivative estimates.
A convenient second-expansion root price is
\[
                \sigma(x)=(1+2^jB_1(x))\omega(x).
\]
Choose a target support exponent $a\ge0$ and a reserve $\eta>0$
with $a+2\eta$ no larger than the available component exponent.
In addition to \eqref{eq:degree-first-smallness}, impose
\begin{equation}\label{eq:degree-second-smallness}
                  \frac{4b_B D\sigma(x)}{\eta}\le\frac12.
\end{equation}
The cluster estimate of Section~\ref{sec:cluster-expansion} now
applies to the strengthened component weights, with root price
$D\sigma$ and tail ratio $1/2$.

For a cluster $\Gamma=(P_1,\ldots,P_q)$, counted with its actual
occurrence multiplicities, define
\[
 d(\Gamma)=\sum_{r=1}^q\{h(P_r)+|P_r|\},\quad
 m(\Gamma)=\sum_r m(P_r),\quad
 s(\Gamma)=\sum_r S(P_r).
\]
Its strengthened weight is its old positive cluster majorant times
\[
             D^{d(\Gamma)}e^{Mm(\Gamma)+\epsilon s(\Gamma)}.
\]
This identity follows by multiplying
\eqref{eq:degree-component-multiplier} over the cluster occurrences;
repeated labels are counted at each occurrence. Suppose
\[
 D^3\le e^M,\qquad D^3\le e^{\epsilon T},\qquad T\ge0.
\]
For any cluster satisfying at least one of
\begin{equation}\label{eq:degree-large-class}
 d(\Gamma)\ge3,\qquad m(\Gamma)\ge1,\qquad s(\Gamma)\ge T,
\end{equation}
the strengthened multiplier is at least $D^3$. Every other factor
is at least one. Its unstrengthened price is therefore bounded by
$D^{-3}$ times the strengthened price.

Spatially weighted derivatives of order at most $j$, with spatial
power $p$, cost $[\ell(x)(1+C(\Gamma))]^{pj}$, where
$\ell(x)=4(2B(x)+1)$. The same exponential-series inequality pays
this polynomial from the reserve $\eta$, at cost
\[
 H_\eta(x)=\ell(x)^{pj}e^\eta(pj)!\eta^{-pj}.
\]
Using the cluster root bound at exponent $a+\eta$, then removing
this reserve, proves the site estimate at exponent $a$
\begin{equation}\label{eq:degree-large-budget}
 \sup_z\sum_{\substack{\Gamma:\ z\in\operatorname{supp}\Gamma\\
                   \Gamma\text{ satisfies }\eqref{eq:degree-large-class}}}
       \|\Phi_\Gamma\|_{j,p}e^{aC(\Gamma)}
                  \le D^{-3}H_\eta(x)\frac{D\sigma(x)}{1-1/2}.
\end{equation}
Here $\Phi_\Gamma$ denotes the cluster contribution and the norm is
its local spatially weighted jet bound. Restricting a nonnegative
sum to the displayed class is legitimate without analyticity of a
field-dependent classification: the class depends only on labels,
supports and charges. Uniform forest averaging preserves the bound.

\subsection{Paying an arbitrary fraction of the forecast}
We check the exponents and the order of choices. Suppose the
unstrengthened site bound is
\begin{equation}\label{eq:degree-old-budget}
                H_i(x)\frac{\sigma(x)}{1-1/2}
                         \le\beta_+(x)^\varphi,
                \qquad \varphi<-3/8,
\end{equation}
for some $i>0$ with $a+2i\le E$, where $E$ is the available component
exponent. This receipt follows from the unstrengthened component
construction when its upper exponent $q$ satisfies $q<\varphi$:
$B_1$ and $H_i$ have upper exponent zero, $\omega$ has upper exponent
$q$, and therefore the entire left side has upper exponent $q$.
Taking the reserve $i=(E-a)/2$ also makes its ordinary degree price
at most $1/2$ eventually, since $q<0$. The same graded local
construction must supply the actual component majorants; the scalar
exponent calculation alone does not construct them.
Put $\eta=(E-a)/2$. Then $0<i\le\eta$ and, writing
$r=pj$, we have
\[
 \eta^{-r}r!\le i^{-r}r!,\qquad
                  H_\eta(x)\le e^\eta H_i(x).
\]
Also $H_i(x)\ge\eta^{-r}r!>0$, since $\ell\ge1$ and $e^i\ge1$.
Thus \eqref{eq:degree-old-budget} bounds $\sigma$ by a fixed constant
times $\beta_+^\varphi$. Since $b_B$ has upper exponent zero and
$3/8+\varphi<0$, it implies
\eqref{eq:degree-second-smallness} at this fixed reserve for all large
$x$. The first condition follows in exactly the same way from the
strictly smaller than $-3/8$ exponent of $r_0$.

The seed mark and the large-footprint threshold obey
$e^{-M}$ and $e^{-\epsilon T}$ estimates with every negative upper
exponent, so the two inequalities $D^3\le e^M,e^{\epsilon T}$ hold
eventually as well. Substituting
\eqref{eq:degree-old-budget} into
\eqref{eq:degree-large-budget} gives
\[
 e^\eta D^{-2}\beta_+^\varphi
                        =e^\eta\beta_+^{\varphi-3/4}.
\]
The incoming regular forecast is $\delta_E(x)=x^{-11/10}$. Since
\[
             \varphi-\frac34+\frac{11}{10}
                         =\varphi+\frac7{20}<0,
\]
the last price is at most $c\delta_E(x)$ for any prescribed $c>0$
after increasing the entrance coupling. The comparison uses
$\beta_+(x)\ge x$ and the negative exponent of $\delta_E$.

For each stage only finitely many transitions and fixed derivative
orders occur. First take the maximum of their block-size requirements;
then take the maximum of their coupling thresholds and the previous
entrance value. This yields one increasing collection of entrance
requirements. Enlarging them preserves all the inequalities above.
It does not require a uniform threshold over infinitely many derivative
orders. The estimates for the remaining degree-one and degree-two
classes, and the identities which attach their subtractions to the
renormalized path, are separate steps.
